\documentclass[12pt,onecolumn]{article}

\usepackage[utf8]{inputenc}
\usepackage[T1]{fontenc}
\usepackage[margin=1in]{geometry}
\usepackage{setspace}
\usepackage{amsmath,amssymb,amsthm}
\usepackage{tikz}
\usetikzlibrary{positioning,fit,backgrounds}

\newcommand\blfootnote[1]{%
\begingroup
\renewcommand\thefootnote{}\footnote{#1}%
\addtocounter{footnote}{-1}%
\endgroup
}

\PassOptionsToPackage{capitalise,nameinlink,noabbrev}{cleveref}
\usepackage[colorlinks=true,linkcolor=blue,citecolor=blue,urlcolor=blue]{hyperref}
\usepackage{cleveref}
\usepackage{orcidlink}
\usepackage{fancyhdr}
\usepackage{xcolor}
\usepackage{graphicx}
\usepackage{mdframed}
\usepackage{xurl}

% Clickable superscript link to definition
\newcommand{\defterm}[2]{#1\textsuperscript{\textcolor{gray}{\hyperref[#2]{\scriptsize$\dagger$}}}}

\crefname{proposition}{Proposition}{Propositions}
\Crefname{proposition}{Proposition}{Propositions}
\crefname{definition}{Definition}{Definitions}
\Crefname{definition}{Definition}{Definitions}
\crefname{remark}{Remark}{Remarks}
\Crefname{remark}{Remark}{Remarks}

\setstretch{1.08}

\pagestyle{fancy}
\fancyhf{}
\setlength{\headheight}{14pt}
\lhead{Time Is Not Global}
\rhead{John C. W. McKinley}
\cfoot{\thepage}

\newcounter{nogo}[section]
\renewcommand{\thenogo}{\thesection.\arabic{nogo}}

\theoremstyle{plain}
\newtheorem{proposition}{Proposition}
\theoremstyle{definition}
\newtheorem{definition}{Definition}
\theoremstyle{remark}
\newtheorem{remark}{Remark}

\makeatletter
\let\c@proposition\c@nogo
\let\c@definition\c@nogo
\let\c@remark\c@nogo
\makeatother

\renewcommand{\theproposition}{\thenogo}
\renewcommand{\thedefinition}{\thenogo}
\renewcommand{\theremark}{\thenogo}

\title{\textbf{Time Is Not Global}\\
\large\textit{Disambiguating the Timeless Light Model from
Wheeler--DeWitt, Page--Wootters, and Relational Quantum Time}}

\author{John C. W. McKinley\,\orcidlink{0009-0005-7097-5035}}

\date{September 2026}

\begin{document}

\maketitle

\blfootnote{\scriptsize Working draft published at \href{https://johncwmckinley.com}{johncwmckinley.com}. This draft has not been deposited on Zenodo and may be revised.}

\begin{abstract}

Several approaches to quantum gravity and quantum foundations remove,
internalize, or reconstruct the ordinary time parameter. The
Wheeler--DeWitt framework does not contain the ordinary external
Schr\"odinger time parameter in its fundamental constraint. Page and
Wootters showed that a globally stationary system can nevertheless exhibit
clock-relative dynamics to an internal observer. Rovelli developed a
framework for observable evolution in a regime in which a fundamental time
variable is not defined. These approaches therefore should not be
characterized as theories that simply retain a universal master time.

The Timeless Light Model (TLM) makes a different claim. Its starting point
is not the absence of a time symbol from an equation. Its starting point is
the architecture of physical law. The laws of the universe determine which
relations and interactions are allowed. They are not events occurring
inside the universe and are not themselves subject to time. A mathematical
expression of a law may contain a temporal parameter where the physical
relation being described includes temporal structure; this does not make
the law itself temporal.

TLM therefore proposes that time is not global. Temporal structure applies
only where the lawful physical relation possesses it. In particular,
special and general relativity govern the temporal and geometric structure
of spacetime registration. A massive timelike system may accumulate proper
time. The photon relation does not: its proper-time interval is null.
Under TLM this is not interpreted as zero elapsed time on an unseen
universal clock. It means that no proper-time interior belongs to the
photon relation itself. Likewise, an entanglement correlation need not be
treated as a third time-bearing process traveling between two local
registrations.

The consequence is a domain restriction on temporal language. A phenomenon
can violate temporal order only if temporal order applies to that relation
in the first place. Expressions such as ``instantaneous,'' ``backward in
time,'' or ``faster than light'' can otherwise import temporal or transit
structure into a lawful relation that does not possess it.

This paper compares TLM with Wheeler--DeWitt, Page--Wootters, and relational
quantum time without treating those approaches as straw men. Considerable
compatibility may exist among them. The distinctive TLM claim is narrower
and architectural: lawful structure is prior to the physical relations it
permits, time is not a universal variable of that structure, and temporal
properties apply only within the lawful relations for which temporal
structure is supplied.

\end{abstract}

\noindent{\small\textcolor{gray}{Technical terms marked with a superscript
dagger (\textsuperscript{\scriptsize$\dagger$}) link to their formal
definition in Section~\ref{sec:definitions}.}}

\begin{center}

\textit{The laws define what may occur.}

\vspace{0.4em}

\textit{The laws themselves do not occur in time.}

\vspace{0.4em}

\textit{Time is not global.}

\vspace{0.4em}

\textit{A temporal violation requires a relation to which time applies.}

\end{center}

\vspace{0.7em}

\clearpage
\tableofcontents
\clearpage


\section{Introduction}

The Timeless Light Model makes a simple claim that can become obscured when
it is placed beside the sophisticated literature on the problem of time.

Time is not global.

This does not mean that physics contains no temporal variables. It does
not mean that clocks are unreal, that duration is unreal, or that equations
containing \(t\) are mistaken. It means that temporal structure should not
be presumed to apply to every lawful physical relation merely because it
applies to some.

The distinction begins with physical law.

The laws of the universe specify the boundaries of allowed physical
relation and interaction. In the terminology of the Timeless Light Model,
the coordinated set of those laws is
\defterm{lawful structure}{def:lawfulstructure}.

The familiar analogy is a lawbook. The analogy should not be taken to mean
that a literal book exists outside the universe. Its purpose is to make the
logical hierarchy clear. The rules determine what physical relations are
permitted. The rules are not themselves physical events occurring under
those same relations.

The laws are therefore architecturally prior to the events and relations
they govern.

``Prior'' here does not mean earlier in time.

There can be no temporal claim that the laws existed for some duration and
then the universe began. Rather, the statement is one of logical
dependence: whatever physical interactions occur must be lawful in order to
occur at all.

Thus,

\[
\boxed{
\text{law}
\longrightarrow
\text{boundary of permitted physical relation}
}
\]

does not mean

\[
\text{law first in time}
\longrightarrow
\text{physical event later in time}.
\]

The distinction matters because a mathematical law may contain a temporal
quantity without the law itself becoming temporal.

A law governing a relation involving \(t\) does not age as \(t\) changes.
It does not wait for the next value of \(t\). It continuously specifies the
allowed relation.

The Timeless Light Model therefore separates two propositions:

\[
\boxed{
\text{a lawful physical relation contains temporal structure}
}
\]

and

\[
\boxed{
\text{law itself is subject to time}.
}
\]

The first may be true.

The second does not follow.

This architecture was developed in the published TLM corpus and stated
directly in \textit{Physics Has Always Been Unified: General Timelessness,
Null Proper Time, and the Clarified Architecture}
\cite{unified2026}.

The purpose of the present paper is narrower. Several important approaches
to quantum gravity also use the language of timelessness. It is therefore
necessary to determine whether TLM merely restates those approaches, differs
from them, or is substantially compatible with them while making an
additional architectural distinction.

That inquiry should begin without caricaturing the neighboring theories.


\section{The Lawbook and the Domain of Time}

\label{sec:lawbook}

The central TLM hierarchy is:

\[
\boxed{\text{LAWFUL STRUCTURE}}
\]

followed not temporally, but logically, by

\[
\boxed{\text{THE PHYSICAL RELATIONS IT PERMITS}.}
\]

Some of those lawful relations possess temporal structure.

Others do not.

This suggests a useful concept: the \emph{domain of time}.

\section{Definitions}
\label{sec:definitions}

\begin{definition}[Lawful structure]
\label{def:lawfulstructure}
Lawful structure is the coordinated specification of the physical laws,
relations, constraints, conservation principles, symmetries, selection
rules, boundary conditions, and coupling structure governing what physical
relations and outcomes are admissible. It is not an event, object, region,
or medium within spacetime.
\end{definition}

\begin{definition}[Architectural priority]
\label{def:architecturalpriority}
Architectural priority is logical dependence rather than temporal
precedence. The laws are architecturally prior to a physical relation in
the sense that the relation must conform to the governing law in order to
be physically admissible. Architectural priority makes no claim that the
laws existed at an earlier time.
\end{definition}

\begin{definition}[Temporal domain]
\label{def:temporaldomain}
The temporal domain is the class of lawful physical relations to which
temporal quantities such as temporal ordering, duration, causal ordering,
or proper-time accumulation apply.
\end{definition}

\begin{definition}[Registration]
\label{def:registration}
Registration is the lawful participation of a physical outcome within
spacetime such that it possesses geometric localization, causal and
temporal relation, and the possibility of observation or interaction by
other registered physical systems.
\end{definition}

\begin{definition}[No-global-time principle]
\label{def:noglobaltime}
The no-global-time principle is the TLM claim that temporal structure is
not presumed to apply to every lawful physical relation. Time applies where
the governing physical relation supplies temporal structure. The existence
of timed relations does not imply a universal temporal variable governing
all lawful structure.
\end{definition}

The last definition is the central claim of this paper.

It can be stated more simply:

\[
\boxed{\text{Time applies where time applies.}}
\]

This statement is intentionally tautological in appearance because it
removes an additional assumption usually inserted without notice: that if
time applies to ordinary spacetime processes, then some form of time must
also govern every physical relation whatsoever.

TLM denies that inference.

A relation possessing a temporal variable is temporal in that respect.

A relation possessing proper time accumulates proper time.

A relation possessing coordinate-time ordering can be described using that
coordinate structure.

But nothing follows about a separate universal variable that must also
govern lawful relations lacking those properties.


\section{A Law Can Contain Time Without Being in Time}

An important source of confusion is the distinction between a law and the
variables appearing in its mathematical representation.

Consider schematically

\[
F(x,t)=0.
\]

The law describes an allowed relationship involving a temporal quantity
\(t\).

Nothing about this notation requires the law itself to pass from \(t_1\)
to \(t_2\).

The law is not younger at \(t_1\) and older at \(t_2\).

It does not develop as the physical system develops.

The relation described by the law may possess temporal structure. The law
specifies the allowed relation.

This distinction permits temporal variables wherever physics requires them
without promoting time into a universal background governing the lawbook
itself.

\begin{proposition}[Law--variable distinction]
\label{prop:lawvariable}
The occurrence of a time variable in the mathematical representation of a
physical law establishes that the represented relation has temporal
structure in the relevant description. It does not establish that the law
itself is subject to time.
\end{proposition}

This proposition is important for comparing TLM with ordinary quantum
mechanics as well as relativity.

The time-dependent Schr\"odinger equation, for example, contains \(t\):

\[
i\hbar\frac{\partial}{\partial t}\Psi
=
\hat H\Psi.
\]

TLM need not remove the variable.

The equation may describe lawful relations involving temporal structure.

The point is only that the law represented by the equation remains the rule
specifying the allowed relation; it is not itself another time-evolving
physical participant.

Special and general relativity make the distinction particularly concrete
because they supply the temporal and geometric organization of registered
spacetime relations.

Thus TLM does not say:

\[
t \text{ is forbidden}.
\]

It says:

\[
\boxed{
t \text{ has a domain of application}.
}
\]


\section{Null Proper Time Is the Diagnostic Case}

The physical case that motivated the TLM reversal is the photon.

For a timelike worldline,

\[
d\tau > 0
\]

and proper time accumulates along the worldline.

For a null worldline,

\[
ds^2 = 0,
\]

and therefore

\[
d\tau = 0.
\]

For light in flat spacetime,

\[
c^2dt^2-dx^2=0.
\]

The usual interpretation notes correctly that a photon has no rest frame
and that proper time along a null worldline is zero.

TLM asks what ontological conclusion should follow.

The conventional traveler picture easily encourages the following
intuition:

\begin{quote}
A universal time still exists everywhere, but the photon's clock somehow
reads zero while the photon travels through it.
\end{quote}

TLM rejects the need for that extra picture.

There is no global temporal variable that must secretly remain attached to
the photon relation.

The emission event registers.

The absorption event registers.

Those endpoints possess coordinate relations within spacetime and are
temporally ordered for registered observers according to relativity.

But the lawful relation called the photon has

\[
\Delta\tau=0.
\]

Under TLM, this means exactly what the physical relation supplies: there is
no accumulated proper-time interior belonging to that relation.

It does not mean

\[
\text{zero units on an otherwise universal hidden clock}.
\]

It means

\[
\boxed{
\text{no proper-time variable is supplied to the photon relation itself}.
}
\]

\begin{proposition}[Null-domain result]
\label{prop:nulldomain}
Null proper time does not require the interpretation that a photon
experiences zero elapsed time relative to a deeper universal time. Under
TLM, the null relation simply lacks an intrinsic proper-time history.
Temporal and geometric ordering belong to the registered endpoints and to
the spacetime description supplied by relativity.
\end{proposition}

This is why null proper time is diagnostic rather than merely numerical.

If time were necessarily global, every physical relation would require a
time-bearing interior even where proper time is exactly zero.

The null case shows that physics itself does not require such universality.


\section{A Temporal Violation Requires a Temporal Domain}

The no-global-time principle produces a simple test.

Before describing a physical phenomenon as violating time, ask:

\begin{quote}
What law assigns temporal structure to the relation alleged to violate it?
\end{quote}

If the relation lies within the temporal domain, then temporal statements
are meaningful.

One can ask whether event \(A\) occurs before event \(B\), whether a signal
travels faster than a permitted propagation speed, or whether a worldline
accumulates a particular proper-time interval.

But if the relation itself contains no corresponding temporal structure,
then applying such language to it imports an additional variable not
supplied by the physical description.

\begin{proposition}[Temporal-violation criterion]
\label{prop:violation}
A lawful relation can violate a temporal constraint only if that temporal
constraint applies to the relation in the first place.
\end{proposition}

This yields the compact TLM statement:

\[
\boxed{
\text{A phenomenon that ``violates time'' may be violating nothing.}
}
\]

The quotation marks are important.

The proposition does not license violations of relativity within the domain
where relativity applies.

A massive object does not evade its timelike worldline.

A controllable signal is not licensed to propagate outside the relativistic
causal structure.

Rather, TLM denies the prior move of assigning a propagation history or
temporal variable to a relation that is not itself a spacetime
transmission.


\section{Retrocausality and the Photon}

Retrocausal interpretations illustrate the importance of the domain test.

To describe a physical influence as retrocausal ordinarily requires a
temporal history within which something associated with a later condition
acts backward toward an earlier condition.

Schematically,

\[
t_2 \longrightarrow t_1,
\qquad t_2>t_1.
\]

For ordinary time-bearing processes, such a proposition is intelligible
whether or not it is correct.

The photon relation presents a different case under TLM.

Its registered endpoints may be ordered:

\[
E_{\mathrm{emit}}
\prec
E_{\mathrm{absorb}}
\]

within the spacetime description.

But the lawful photon relation itself has no accumulated proper-time
interior.

Assigning to that relation a forward internal history and then reversing
part of that history therefore introduces precisely the structure that the
null case does not supply.

Under TLM, the issue is not that the photon travels forward rather than
backward through a hidden time.

The claim is:

\[
\boxed{
\text{there is no photon proper-time history available to reverse}.
}
\]

The retrocausal problem is therefore dissolved rather than solved by a
different causal mechanism \cite{retrocausality}.

This does not erase the coordinate ordering of emission and absorption.

It rejects the additional traveler history inserted between them.


\section{Entanglement and the Absence of Temporal Transit}

Quantum entanglement provides a second application.

An entangled system produces correlated local outcomes.

Each detector event is a registration.

Each detector has a spacetime location.

Each registered outcome participates in temporal and causal relations.

Relativity applies to those registrations and to any ordinary physical
signal exchanged between them.

The entanglement correlation is different in kind.

Under TLM, the joint correlation belongs to the lawful structure governing
the pair.

It is not a third detector event.

It is not an object traveling from detector \(A\) to detector \(B\).

It is not a signal whose speed must be calculated by

\[
v=\frac{\Delta x}{\Delta t}.
\]

Consequently, describing the correlation as ``instantaneous'' already
risks assigning a temporal transit variable to something that is not a
transit.

The architecture is instead:

\[
\boxed{
\text{one lawful relation}
}
\]

with

\[
\boxed{
\text{two local spacetime registrations}.
}
\]

The registrations are timed.

The lawful correlation need not be.

\begin{proposition}[Entanglement-domain result]
\label{prop:entanglement}
Under TLM, the temporal and causal constraints of relativity govern the
local registrations of an entangled system and any physical communication
between them. The joint correlation is not thereby required to possess a
propagation time between the registrations.
\end{proposition}

Thus the familiar question

\begin{quote}
How does the correlation travel between the detectors so quickly?
\end{quote}

contains a TLM category error.

It assumes that there is a traveling correlation-bearing process to which
speed and duration apply.

If there is no such process, there is no superluminal temporal violation to
explain.

There is a lawful joint relation and there are relativistically organized
local registrations \cite{unified2026}.


\section{What Wheeler--DeWitt Actually Says}

The preceding architecture should now be compared with the quantum-gravity
literature without simplifying that literature into a foil.

In canonical quantum gravity, the Wheeler--DeWitt constraint is written
schematically as

\[
\hat{\mathcal H}\Psi=0.
\]

Unlike the ordinary Schr\"odinger equation, this expression contains no
external time derivative of the familiar form

\[
i\hbar\frac{\partial}{\partial t}\Psi.
\]

DeWitt's canonical treatment makes the point more carefully than the common
popular description of a ``timeless equation.'' For a finite world, the
constraints govern the dynamics; the coordinate label \(x^0\) is irrelevant,
and ``time'' must be determined intrinsically \cite{dewitt1967}.

This already rejects a simple external universal clock.

TLM should therefore not claim that Wheeler--DeWitt retains one.

The difference lies elsewhere.

The Wheeler--DeWitt problem is a problem within a physical formalism:

\begin{quote}
How is dynamical physics represented when ordinary external Schr\"odinger
time is absent?
\end{quote}

The TLM question is architectural:

\begin{quote}
Which lawful relations possess temporal structure at all?
\end{quote}

The absence of \(t\) from an equation is neither necessary nor sufficient
for the TLM conclusion.

A law may contain \(t\) while the law itself remains timeless.

A law may omit \(t\) without thereby proving anything about the global
ontology of time.

Thus,

\[
\boxed{
\text{formal absence of external }t
\neq
\text{TLM no-global-time principle}.
}
\]

The two claims may be compatible without being identical.


\section{Page--Wootters: Clock Relations Without External Time}

Page and Wootters begin even closer to the issue.

They observe that the time parameter in the Schr\"odinger equation is not
itself observable and show that it is consistent with observation to treat
a closed system such as the universe as stationary. Observed dynamical
evolution can then be described through dependence upon internal clock
readings \cite{pagewootters1983}.

Schematically, let the total system contain a clock subsystem \(C\) and
another subsystem \(S\):

\[
\mathcal H
=
\mathcal H_C\otimes\mathcal H_S.
\]

The internal observer obtains relational structure such as

\[
C_i\leftrightarrow S_i.
\]

Page--Wootters therefore should not be presented as a theory in which a
master external time continues to control the universe.

It explicitly shows how ordinary dynamics can be represented without one.

This creates a genuine point of contact with TLM.

The TLM question becomes:

\begin{quote}
What ontological conclusion should be drawn from
\(C_i\leftrightarrow S_i\)?
\end{quote}

One possible description is that time has emerged.

TLM permits a more economical description:

\[
\boxed{
\text{a clock-relative temporal relation has been established}.
}
\]

The clock states and the compared system states provide registered physical
relations.

That is sufficient for ordering and temporal description.

Nothing further requires lawful structure itself to acquire a global time
variable.

\begin{remark}[Compatibility with Page--Wootters]
The present paper does not claim that Page--Wootters is incompatible with
TLM. To the extent that Page--Wootters eliminates external time and treats
observable dynamics relationally, the frameworks may be strongly
compatible. The TLM contribution is the additional architectural claim
that the clock-relative temporal structure should not automatically be
promoted into a universal property of all lawful relations.
\end{remark}


\section{Moreva et al.: What the Experiment Shows}

Moreva et al.\ experimentally illustrated the Page--Wootters mechanism
using entangled photon polarizations \cite{moreva2014}.

Their description is particularly useful because it cleanly distinguishes
two viewpoints.

An internal observer correlated with the clock subsystem sees the other
subsystem evolve.

An external observer considering the appropriate global properties can
show that the joint state is static.

The experiment therefore illustrates that clock-relative dynamics can
exist within a globally stationary description.

TLM need not dispute any part of this result.

The question is instead how to name what has been established.

The physical observations include:

\[
\text{clock states},
\]

\[
\text{system states},
\]

and

\[
\text{correlations between them}.
\]

Those relations can supply an operational temporal description.

Calling that description ``emergent time'' is unobjectionable if
\emph{emergent time} means only that temporal ordering becomes available
to an internal observer through physical clock relations.

A stronger conclusion would be that a global physical constituent called
time has thereby been created.

That stronger conclusion is not required by the experiment itself.

The TLM description is simply:

\[
\boxed{
\text{clock-relative temporal structure appears where the physical relation
supplies it}.
}
\]


\section{Rovelli: Physics Without a Fundamental Time Variable}

Rovelli comes still closer to the no-global-time question.

He explicitly considers the hypothesis that time is not defined at the
fundamental level and develops gauge-invariant observables capable of
describing observable evolution in that regime \cite{rovelli1991}. Earlier
work likewise studied quantum mechanics in a regime where unitary time
evolution appears only as an approximation \cite{rovelli1990}.

It would therefore be particularly misleading to characterize Rovelli as
assuming that a universal time variable remains present underneath the
theory.

He explicitly investigates the opposite possibility.

The proper comparison is consequently not

\[
\text{Rovelli has global time}
\qquad\text{versus}\qquad
\text{TLM does not}.
\]

It is:

\[
\boxed{
\text{What does ``no fundamental time'' mean in each architecture?}
}
\]

Rovelli develops a formalism for observable physics in the absence of a
fundamental time variable.

TLM develops an ontological classification of physical law and physical
relation in which temporal structure has a restricted domain.

These may be complementary statements.

TLM adds the specific claim that relations such as the null photon relation
and the entanglement correlation should not be assigned a temporal history
merely because their registered endpoints or outcomes participate in
spacetime.


\section{The Difference Between Emergent Time and Domain-Limited Time}

The comparison suggests that the phrase \emph{emergent time} can obscure
two different ideas.

The first is operational.

A subsystem can serve as a clock.

Correlated states can be ordered.

An internal observer can describe evolution.

A macroscopic description can contain durations and histories.

TLM has no objection to calling this operational temporal emergence.

The second idea is ontological:

\begin{quote}
A fundamentally timeless reality must somehow generate a new global
physical factor called time.
\end{quote}

TLM does not require this step.

The TLM architecture is not

\[
\text{timeless reality}
\longrightarrow
\text{time comes into existence}
\longrightarrow
\text{temporal universe}.
\]

Such a picture is already awkward because the arrow appears to describe a
temporal transition occurring before time exists.

The TLM architecture is instead domain-based:

\[
\boxed{
\text{lawful structure}
}
\]

contains lawful relations of different kinds.

Some possess temporal structure.

Some do not.

Where spacetime registration occurs, temporal and geometric relations are
governed by special and general relativity.

Where a lawful relation contains no internal temporal variable, TLM does
not add one.

There is therefore no universal transition from timelessness into time.

There is only a distinction between relations to which temporal structure
applies and relations to which it does not.

\section{The Central Comparison}

The neighboring approaches can now be compared without assigning claims to
them that they do not make.

\begin{center}
\begin{tabular}{p{0.20\linewidth}p{0.35\linewidth}p{0.34\linewidth}}
\hline
\textbf{Framework} &
\textbf{Central time question} &
\textbf{Relation to TLM}
\\
\hline

Wheeler--DeWitt &
How is quantum-gravitational dynamics represented when ordinary external
Schr\"odinger time is absent? &
Compatible with no external master clock. TLM additionally asks which
lawful relations possess temporal structure at all.
\\[0.8em]

Page--Wootters &
How can observed dynamics be described relationally using an internal
physical clock while the closed global system is stationary? &
Strongly compatible with relational temporal description. TLM does not
infer from clock relations that time must be global.
\\[0.8em]

Moreva et al. &
Can Page--Wootters clock-relative dynamics be experimentally illustrated? &
The observed correlations demonstrate clock-relative temporal structure;
they need not establish a new universal time constituent.
\\[0.8em]

Rovelli &
How can observable evolution be described when a fundamental time variable
is not defined? &
Close conceptual neighbor. TLM adds a domain-based ontology connecting
no-global-time to null proper time, photon interpretation, entanglement,
and registration.
\\[0.8em]

Timeless Light Model &
Which lawful relations possess temporal structure, and which do not? &
Time is not global. Laws define the permitted relations; temporal structure
applies only where the physical relation supplies it.
\\

\hline
\end{tabular}
\end{center}


\section{What TLM Adds}

If Wheeler--DeWitt, Page--Wootters, and Rovelli already remove external or
fundamental time in important respects, what remains distinctive about TLM?

The answer is not that TLM discovered mathematical timelessness first.

It did not.

Nor is the answer that existing approaches secretly retain universal time.

They need not.

The TLM contribution proposed here consists of three connected claims.

First, lawful structure is architecturally senior to the physical relations
it permits. Laws determine the boundaries of admissible interaction and
relation but are not themselves events passing through time.

Second, temporal structure has a domain rather than universal scope.

\[
\boxed{
\text{Time is not global.}
}
\]

Third, this domain restriction is applied directly to physical puzzles
normally formulated using global temporal language.

For the photon:

\[
\Delta\tau=0
\]

is read as the absence of a proper-time interior rather than as zero
progress on an unseen universal clock.

For entanglement:

\[
\text{one lawful relation}
+
\text{two timed local registrations}
\]

does not require a time-bearing signal crossing between the detectors.

For retrocausality:

\[
\text{no internal photon time history}
\]

means that there is no such history to reverse.

These are not three independent interpretive tricks.

They are applications of one rule:

\[
\boxed{
\text{Do not assign time to a lawful relation unless its governing
physical structure supplies time.}
}
\]


\section{What This Paper Does Not Claim}

The distinction is intentionally narrow.

\begin{remark}
This paper does not claim that Wheeler--DeWitt, Page--Wootters, or Rovelli
assume an external universal master clock. Their published work explicitly
moves away from such a description.
\end{remark}

\begin{remark}
This paper does not claim that equations containing a time parameter are
incorrect or less fundamental merely because the parameter appears.
\end{remark}

\begin{remark}
This paper does not claim that clocks, duration, temporal order, causal
structure, or proper time are unreal.
\end{remark}

\begin{remark}
This paper does not claim that the registered endpoints of a null relation
lack coordinate-time ordering for ordinary spacetime observers.
\end{remark}

\begin{remark}
This paper does not permit controllable signals to violate the relativistic
causal structure. Its claim concerning entanglement is that the correlation
itself should not first be assumed to be a propagating signal.
\end{remark}

\begin{remark}
This paper does not provide a replacement quantum-gravity formalism.
\end{remark}

\begin{remark}
This paper does not claim that all uses of the phrase ``emergent time'' are
incompatible with TLM. If the phrase means the appearance of operational
clock-relative temporal description, the two may be entirely compatible.
\end{remark}


\section{Discussion}

The problem of time may contain more than one problem.

There is a formal problem concerning the representation of dynamics when
the ordinary external time parameter disappears.

There is an operational problem concerning how internal observers construct
clocks and temporal descriptions.

There is a relational problem concerning how one physical quantity changes
with respect to another.

And there is an ontological question:

\begin{quote}
Must time be assigned globally to every lawful physical relation?
\end{quote}

TLM answers no.

This answer is intentionally simple.

The laws specify what may happen.

They do not themselves happen in time.

Some of the relations permitted by those laws participate within spacetime
and possess temporal structure.

Others do not contain an internal temporal variable.

There is no requirement that a second hidden time continue running behind
the relations for which the physical description supplies none.

This is why null proper time matters so strongly.

The equation

\[
\Delta\tau=0
\]

does not need to be translated into

\begin{quote}
zero elapsed units of a deeper time that still universally exists.
\end{quote}

It can be read literally as a limit on the temporal structure assigned to
the null relation.

The same discipline clarifies entanglement.

The two registrations have temporal coordinates.

The lawful correlation need not have a transit time.

And it clarifies retrocausality.

The registered endpoints have temporal order.

The photon relation need not possess an intermediate proper-time history
that can be made to run backward.

Thus the TLM maxim can be stated:

\[
\boxed{
\text{Temporal predicates apply only inside the temporal domain.}
}
\]

The statement sounds modest.

Its consequences are not.

If correct, several phenomena described as challenging temporal order may
instead expose the mistaken assumption that temporal order was universal.


\section{Conclusion}

The Timeless Light Model and the existing problem-of-time literature share
substantial conceptual territory.

Wheeler--DeWitt removes the ordinary external Schr\"odinger time parameter
from the canonical quantum-gravity constraint and requires time to be
treated intrinsically.

Page and Wootters show that a closed stationary system can nevertheless
support observed dynamics through internal clock readings.

Moreva et al.\ experimentally illustrate such clock-relative dynamics.

Rovelli explicitly develops quantum description in a regime where a
fundamental time variable is not defined.

TLM need not oppose any of these results.

Its claim lies at another level.

The laws of the universe define the boundaries of permitted interaction
and relation.

The laws are not themselves events.

They do not wait, age, or evolve.

A law may describe a relation containing temporal variables without the law
itself becoming subject to time.

Most importantly,

\[
\boxed{
\text{time is not global}.
}
\]

Temporal structure applies where the governing physical relation supplies
it.

For timelike spacetime relations, proper time applies.

For registered spacetime events, temporal and causal ordering apply.

For the null photon relation,

\[
\Delta\tau=0,
\]

and TLM adds no hidden proper-time history behind that result.

For an entanglement correlation, the local outcomes register within
spacetime, but the correlation need not be assigned a propagation time
between them.

The resulting principle is therefore:

\[
\boxed{
\begin{minipage}{0.82\linewidth}
A phenomenon can violate time only where time applies. If a lawful relation
does not possess the relevant temporal structure, descriptions such as
instantaneous, retrocausal, or faster than light may import a variable that
the relation itself does not contain.
\end{minipage}
}
\]

Timelessness in TLM is therefore not the claim that the entire universe is
frozen at zero time.

It is not the claim that temporal variables should be removed from physics.

It is not the claim that clocks are illusions.

It is the simpler claim that the universe does not possess one global time
variable that must govern every lawful relation.

The laws define what is allowed.

Where those laws supply time, time applies.

Where they do not, nothing is missing.


\begin{thebibliography}{99}

\bibitem{unified2026}
J.~C.~W.~McKinley,
\textit{Physics Has Always Been Unified:
General Timelessness, Null Proper Time, and the Clarified Architecture ---
Why Quantum Mechanics and Relativity Were Never in Competition}.
Zenodo, 2026.
\href{https://doi.org/10.5281/zenodo.20954931}
{doi:10.5281/zenodo.20954931}.

\bibitem{bedrock}
J.~C.~W.~McKinley,
\textit{A Minimal Structural Statement of the Timeless Light Model}.
Zenodo, 2026.
\href{https://doi.org/10.5281/zenodo.18521383}
{doi:10.5281/zenodo.18521383}.

\bibitem{generaltimelessness}
J.~C.~W.~McKinley,
\textit{Timelessness Is the General Condition:
Spacetime Is the Special Case}.
Zenodo, 2026.
\href{https://doi.org/10.5281/zenodo.19771924}
{doi:10.5281/zenodo.19771924}.

\bibitem{nullpropertime}
J.~C.~W.~McKinley,
\textit{Taking Null Proper Time Seriously:
An Interpretive Clarification of Null Proper Time}.
Zenodo, 2025.
\href{https://doi.org/10.5281/zenodo.18004632}
{doi:10.5281/zenodo.18004632}.

\bibitem{retrocausality}
J.~C.~W.~McKinley,
\textit{Retrocausal Objections Are Disallowed for the Photon}.
Zenodo, 2026.
\href{https://doi.org/10.5281/zenodo.19648208}
{doi:10.5281/zenodo.19648208}.

\bibitem{dewitt1967}
B.~S.~DeWitt,
\textit{Quantum Theory of Gravity. I. The Canonical Theory}.
Physical Review \textbf{160}, 1113--1148, 1967.
\href{https://doi.org/10.1103/PhysRev.160.1113}
{doi:10.1103/PhysRev.160.1113}.

\bibitem{pagewootters1983}
D.~N.~Page and W.~K.~Wootters,
\textit{Evolution without evolution:
Dynamics described by stationary observables}.
Physical Review D \textbf{27}, 2885--2892, 1983.
\href{https://doi.org/10.1103/PhysRevD.27.2885}
{doi:10.1103/PhysRevD.27.2885}.

\bibitem{moreva2014}
E.~Moreva, G.~Brida, M.~Gramegna, V.~Giovannetti,
L.~Maccone, and M.~Genovese,
\textit{Time from quantum entanglement:
An experimental illustration}.
Physical Review A \textbf{89}, 052122, 2014.
\href{https://doi.org/10.1103/PhysRevA.89.052122}
{doi:10.1103/PhysRevA.89.052122}.

\bibitem{rovelli1990}
C.~Rovelli,
\textit{Quantum mechanics without time: A model}.
Physical Review D \textbf{42}, 2638--2646, 1990.
\href{https://doi.org/10.1103/PhysRevD.42.2638}
{doi:10.1103/PhysRevD.42.2638}.

\bibitem{rovelli1991}
C.~Rovelli,
\textit{Time in quantum gravity: An hypothesis}.
Physical Review D \textbf{43}, 442--456, 1991.
\href{https://doi.org/10.1103/PhysRevD.43.442}
{doi:10.1103/PhysRevD.43.442}.

\end{thebibliography}

\end{document}